% Generated from locale/tr/content/normal-modal-logic/tableaux/countermodels.tex; do not hand-edit.


\olfileid[tr]{nml}{tab}{cou}
      
\olsection{\usetoken{P}{tableau}dan Karşı-Modeller}

\begin{explain}
  Tamlık teoreminin ispatı yalnızca $\Entails !A$ ise $\Proves !A$ olduğunu
  göstermekle kalmaz; $\Entails/ !A$ ise $!A$ için karşı-model kurma yöntemi
  de verir. $\Log{K}$ durumunda bu yöntem bir \emph{karar yordamı} oluşturur.
  Nitekim $\Entails/ !A$ olsun. \olref[cpl]{prop:complete-tableau}
  önermesinin ispatı, tam bir !!{tableau} kurma yöntemi verir. Bu yöntem
  gerçekte her zaman sona erer. $\Log{K}$ için önermesel kurallar yalnızca daha
  düşük karmaşıklıkta önekli !!{formula}s ekler; yani her önermesel kuralın,
  herhangi bir $\sFmla{S}{!A}[\sigma]$ işaretli formülü için bir dalda yalnızca
  bir kez uygulanması gerekir. Yeni önekleri yalnızca
  \iftag{prvBox}{$\TRule{\False}{\Box}$}{}\iftag{notprvBox,notprvDiamond}{}{
    ve }\iftag{prvDiamond}{$\TRule{\True}{\Diamond}$}{}
  \iftag{notprvBox,notprvDiamond}{kuralı}{kuralları} üretir; bunların da
  yalnızca bir kez uygulanması gerekir (ve tek bir yeni önek üretirler).
  \iftag{prvBox}{$\TRule{\True}{\Box}$}{}\iftag{notprvBox,notprvDiamond}{}{
    ve }\iftag{prvDiamond}{$\TRule{\False}{\Diamond}$}{}
  \iftag{notprvBox,notprvDiamond}{kuralının}{kurallarının} birden çok kez
  uygulanması gerekebilir; fakat her önek için yalnızca bir kez uygulanırlar
  ve yalnızca sonlu sayıda yeni önek üretilir. Dolayısıyla kuruluş, sonlu sayıda
  aşamadan sonra ya kapalı bir dal ya da tam bir dal verir.

  Açık ve tam bir dalı bulunan tablo kurulduğunda,
  \olref[cpl]{thm:tableau-completeness} teoreminin ispatı özgün önekli
  !!{formula}s kümesini sağlayan açık bir model verir. Dolayısıyla
  $\Gamma \Entails !A$ ise kapalı bir !!{tableau} bulunması ve
  $\Gamma \Proves !A$ olması bir yana, kapalı !!{tableau}{}yu doğru biçimde
  arayıp ``tam'' bir !!{tableau} elde edersek yalnızca
  $\Gamma \Entails/ !A$ olduğunu bilmekle kalmaz, gerçekten bir karşı-model de
  kurabiliriz.
\end{explain}

\iftag{prvBox}{
\begin{ex}
  $\Proves/ \Box(p \lor q) \lif (\Box p \lor \Box q)$ olduğunu biliyoruz.
  Bir tablonun kurulması şöyle başlar:
  \begin{oltableau}
    [\pFmla{\False}{\Box(p \lor q) \lif (\Box p \lor \Box q)}{1},
      just = \TAss, checked
      [\pFmla{\True}{\Box(p \lor q)}{1},
        just = {\TRule{\False}{\lif}[1]}, 
        [\pFmla{\False}{\Box p \lor \Box q}{1},
          just = {\TRule{\False}{\lif}[1]}, checked
          [\pFmla{\False}{\Box p}{1},
            just = {\TRule{\False}{\lor}[3]}, checked
            [\pFmla{\False}{\Box q}{1},
              just = {\TRule{\False}{\lor}[3]}, checked
              [\pFmla{\False}{p}{1.1}, 
                just = {\TRule{\False}{\Box}[4]}, checked
                [\pFmla{\False}{q}{1.2}, 
                  just = {\TRule{\False}{\Box}[5]}, checked
                ]
              ]
            ]
          ]
        ]
      ]
    ]
  \end{oltableau}
  !!{tableau} elbette henüz tamamlanmamıştır. Sonraki adımda, onay işareti
  olmayan tek satırı, yani $2$.~satırdaki önekli !!{formula}
  $\sFmla{\True}{\Box(p \lor q)}[1]$'i ele alırız. Kapalı tablonun kuruluşu,
  dalda kullanılan her önek, yani hem $1.1$ hem $1.2$ için
  $\TRule{\True}{\Box}$ kuralını uygulamamızı söyler:
  \begin{oltableau}
    [\pFmla{\False}{\Box(p \lor q) \lif (\Box p \lor \Box q)}{1},
      just = \TAss, checked
      [\pFmla{\True}{\Box(p \lor q)}{1},
        just = {\TRule{\False}{\lif}[1]}, 
        [\pFmla{\False}{\Box p \lor \Box q}{1},
          just = {\TRule{\False}{\lif}[1]}, checked
          [\pFmla{\False}{\Box p}{1},
            just = {\TRule{\False}{\lor}[3]}, checked
            [\pFmla{\False}{\Box q}{1},
              just = {\TRule{\False}{\lor}[3]}, checked
              [\pFmla{\False}{p}{1.1}, 
                just = {\TRule{\False}{\Box}[4]}, checked
                [\pFmla{\False}{q}{1.2}, 
                  just = {\TRule{\False}{\Box}[5]}, checked
                  [\pFmla{\True}{p \lor q}{1.1}, 
                    just = {\TRule{\True}{\Box}[2]}
                    [\pFmla{\True}{p \lor q}{1.2}, 
                      just = {\TRule{\True}{\Box}[2]}
                    ]
                  ]
                ]
              ]
            ]
          ]
        ]
      ]
    ]
  \end{oltableau}
  Artık 2., 8. ve 9.~satırlarda onay işareti yoktur. Fakat yeni önek
  eklenmediğinden, kapanmadıkları sürece ortaya çıkan bütün dallarda 8. ve
  9.~satırlara $\TRule{\True}{\lor}$ uygularız:
  \begin{oltableau}
    [\pFmla{\False}{\Box(p \lor q) \lif (\Box p \lor \Box q)}{1},
      just = \TAss, checked
      [\pFmla{\True}{\Box(p \lor q)}{1},
        just = {\TRule{\False}{\lif}[1]}, checked
        [\pFmla{\False}{\Box p \lor \Box q}{1},
          just = {\TRule{\False}{\lif}[1]}, checked
          [\pFmla{\False}{\Box p}{1},
            just = {\TRule{\False}{\lor}[3]}, checked
            [\pFmla{\False}{\Box q}{1},
              just = {\TRule{\False}{\lor}[3]}, checked
              [\pFmla{\False}{p}{1.1}, 
                just = {\TRule{\False}{\Box}[4]}, checked
                [\pFmla{\False}{q}{1.2}, 
                  just = {\TRule{\False}{\Box}[5]}, checked
                  [\pFmla{\True}{p \lor q}{1.1}, 
                    just = {\TRule{\True}{\Box}[2]}, checked
                    [\pFmla{\True}{p \lor q}{1.2}, 
                      just = {\TRule{\True}{\Box}[2]}, checked
                      [\pFmla{\True}{p}{1.1},
                        just = {\TRule{\True}{\lor}[8]}, checked, close
                      ]
                      [\pFmla{\True}{q}{1.1},
                        just = {\TRule{\True}{\lor}[8]}, checked
                        [\pFmla{\True}{p}{1.2},
                          just = {\TRule{\True}{\lor}[9]}, checked]
                        [\pFmla{\True}{q}{1.2},
                          just = {\TRule{\True}{\lor}[9]}, checked, close]
                      ]
                    ]
                  ]
                ]
              ]
            ]
          ]
        ]
      ]
    ]
  \end{oltableau}
  Geriye bir açık dal kalır ve bu dal tamdır. Ondan, dünyaları
  $W = \{1, 1.1, 1.2\}$ (açık dalda geçen yegâne önekler), erişilebilirlik
  bağıntısı $R = \{\tuple{1, 1.1}, \tuple{1, 1.2}\}$ ve değerlemesi
  $V(p)=\{1.2\}$ (çünkü 11.~satır $\sFmla{\True}{p}[1.2]$ içerir) ile
  $V(q)=\{1.1\}$ (çünkü 10.~satır $\sFmla{\True}{q}[1.1]$ içerir) olan modeli
  tanımlarız. Model \olref{fig:counter-Box}'ta gösterilmiştir ve bunun
  $\Box(p \lor q) \lif (\Box p \lor \Box q)$ için bir karşı-model olduğunu
  doğrulayabilirsiniz.
  \begin{figure}
  \begin{center}
    \begin{tikzpicture}[modal]
      \node[world] (w1) [label={[align=right]right:\mFalse{p}\\ \mFalse{q}}]{$1$}; 
      \node[world] (w2) [label={[align=right]right:\mFalse{p}\\ \mTrue{q}},
        above left=of w1]{$1.1$}; 
      \node[world] (w3) [label={[align=right]right:\mTrue{p}\\ \mFalse{q}},
        above right=of w1] {$1.2$};
      \draw[->] (w1) to (w2);
      \draw[->] (w1) to (w3);
    \end{tikzpicture}
  \end{center}
  \caption{$\Box(p \lor q) \lif (\Box p \lor \Box q)$ için bir karşı-model.}
  \ollabel{fig:counter-Box}
\end{figure}
\end{ex}          
}
{
\begin{ex}
  $\Proves/ (\Diamond p \land \Diamond q) \lif \Diamond(p \land q)$
  olduğunu biliyoruz. Bir tablonun kurulması şöyle başlar:
  \begin{oltableau}
    [\pFmla{\False}{(\Diamond p \land \Diamond q) \lif \Diamond(p \land q)}{1},
      just = \TAss, checked
      [\pFmla{\True}{\Diamond p \land \Diamond q}{1},
        just = {\TRule{\False}{\lif}[1]}, checked
        [\pFmla{\False}{\Diamond(p \land q)}{1},
          just = {\TRule{\False}{\lif}[1]}
          [\pFmla{\True}{\Diamond p}{1},
            just = {\TRule{\True}{\land}[2]}, checked
            [\pFmla{\True}{\Diamond q}{1},
              just = {\TRule{\True}{\land}[2]}, checked
              [\pFmla{\True}{p}{1.1}, 
                just = {\TRule{\True}{\Diamond}[4]}, checked
                [\pFmla{\True}{q}{1.2}, 
                  just = {\TRule{\True}{\Diamond}[5]}, checked
                ]
              ]
            ]
          ]
        ]
      ]
    ]
  \end{oltableau}
  !!{tableau} elbette henüz tamamlanmamıştır. Sonraki adımda, onay işareti
  olmayan tek satırı, yani $3$.~satırdaki önekli !!{formula}
  $\sFmla{\False}{\Diamond(p \land q)}[1]$'i ele alırız. Kapalı tablonun
  kuruluşu, dalda kullanılan her önek, yani hem $1.1$ hem $1.2$ için
  $\TRule{\False}{\Diamond}$ kuralını uygulamamızı söyler:
  \begin{oltableau}
    [\pFmla{\False}{(\Diamond p \land \Diamond q) \lif \Diamond(p \land q)}{1},
      just = \TAss, checked
      [\pFmla{\True}{\Diamond p \land \Diamond q}{1},
        just = {\TRule{\False}{\lif}[1]}, checked
        [\pFmla{\False}{\Diamond (p \land q)}{1},
          just = {\TRule{\False}{\lif}[1]}
          [\pFmla{\True}{\Diamond p}{1},
            just = {\TRule{\True}{\land}[2]}, checked
            [\pFmla{\True}{\Diamond q}{1},
              just = {\TRule{\True}{\land}[2]}, checked
              [\pFmla{\True}{p}{1.1}, 
                just = {\TRule{\True}{\Diamond}[4]}, checked
                [\pFmla{\True}{q}{1.2}, 
                  just = {\TRule{\True}{\Diamond}[5]}, checked
                  [\pFmla{\False}{p \land q}{1.1}, 
                    just = {\TRule{\False}{\Diamond}[3]}
                    [\pFmla{\False}{p \land q}{1.2}, 
                      just = {\TRule{\False}{\Diamond}[3]}
                    ]
                  ]
                ]
              ]
            ]
          ]
        ]
      ]
    ]
  \end{oltableau}
  Artık 3., 8. ve 9.~satırlarda onay işareti yoktur. Fakat yeni önek
  eklenmediğinden, kapanmadıkları sürece ortaya çıkan bütün dallarda 8. ve
  9.~satırlara $\TRule{\False}{\land}$ uygularız:
  \begin{oltableau}
    [\pFmla{\False}{(\Diamond p \land \Diamond q) \lif \Diamond(p \land q)}{1},
      just = \TAss, checked
      [\pFmla{\True}{\Diamond p \land \Diamond q}{1},
        just = {\TRule{\False}{\lif}[1]}, checked
        [\pFmla{\False}{\Diamond(p \land q)}{1},
          just = {\TRule{\False}{\lif}[1]}
          [\pFmla{\True}{\Diamond p}{1},
            just = {\TRule{\True}{\land}[2]}, checked
            [\pFmla{\True}{\Diamond q}{1},
              just = {\TRule{\True}{\land}[2]}, checked
              [\pFmla{\True}{p}{1.1}, 
                just = {\TRule{\True}{\Diamond}[4]}, checked
                [\pFmla{\True}{q}{1.2}, 
                  just = {\TRule{\True}{\Diamond}[5]}, checked
                  [\pFmla{\False}{p \land q}{1.1}, 
                    just = {\TRule{\False}{\Diamond}[3]}, checked
                    [\pFmla{\False}{p \land q}{1.2}, 
                      just = {\TRule{\False}{\Diamond}[3]}, checked
                      [\pFmla{\False}{p}{1.1},
                        just = {\TRule{\False}{\land}[8]}, checked, close
                      ]
                      [\pFmla{\False}{q}{1.1},
                        just = {\TRule{\False}{\land}[8]}, checked
                        [\pFmla{\False}{p}{1.2},
                          just = {\TRule{\False}{\land}[9]}, checked]
                        [\pFmla{\False}{q}{1.2},
                          just = {\TRule{\False}{\land}[9]}, checked, close]
                      ]
                    ]
                  ]
                ]
              ]
            ]
          ]
        ]
      ]
    ]
  \end{oltableau}
  Geriye bir açık dal kalır ve bu dal tamdır. Ondan, dünyaları
  $W = \{1, 1.1, 1.2\}$ (açık dalda geçen yegâne önekler), erişilebilirlik
  bağıntısı $R = \{\tuple{1, 1.1}, \tuple{1, 1.2}\}$ ve değerlemesi
  $V(p)=\{1.1\}$ (çünkü 6.~satır $\sFmla{\True}{p}[1.1]$ içerir) ile
  $V(q)=\{1.2\}$ (çünkü 7.~satır $\sFmla{\True}{q}[1.2]$ içerir) olan modeli
  tanımlarız. Model \olref{fig:counter-Diamond}'da gösterilmiştir ve bunun
  $(\Diamond p \land \Diamond q) \lif \Diamond (p \land q)$ için bir
  karşı-model olduğunu doğrulayabilirsiniz.
  \begin{figure}
  \begin{center}
    \begin{tikzpicture}[modal]
      \node[world] (w1) [label={[align=right]right:\mFalse{p}\\ \mFalse{q}}]{$1$}; 
      \node[world] (w2) [label={[align=right]right:\mTrue{p}\\ \mFalse{q}},
        above left=of w1]{$1.1$}; 
      \node[world] (w3) [label={[align=right]right:\mFalse{p}\\ \mTrue{q}},
        above right=of w1] {$1.2$};
      \draw[->] (w1) to (w2);
      \draw[->] (w1) to (w3);
    \end{tikzpicture}
  \end{center}
  \caption{$(\Diamond p \land \Diamond q) \lif
    \Diamond (p \land q)$ için bir karşı-model.}
  \ollabel{fig:counter-Diamond}
\end{figure}
\end{ex}          
}

